\documentclass[11pt,a4paper]{article}
\usepackage[utf8]{inputenc}
\usepackage{amsmath,amssymb}
\usepackage{geometry}
\geometry{margin=2.5cm}
\usepackage{listings}
\usepackage{xcolor}
\usepackage{hyperref}

\lstset{
  language=C++,
  basicstyle=\ttfamily\small,
  keywordstyle=\color{blue}\bfseries,
  commentstyle=\color{green!60!black},
  stringstyle=\color{red!70!black},
  numbers=left,
  numberstyle=\tiny\color{gray},
  numbersep=8pt,
  frame=single,
  breaklines=true,
  tabsize=4,
  showstringspaces=false
}

\title{IOI 1994 -- The Triangle}
\author{Solution}
\date{}

\begin{document}
\maketitle

\section{Problem Statement}

Given a triangle of $n$ rows of positive integers, find a path from the top to the bottom
such that the sum of the numbers along the path is maximized.
At each step, you may move to the element directly below or to the element
directly below and one position to the right.

\medskip
\noindent\textbf{Example:}
\begin{verbatim}
        7
       3 8
      8 1 0
     2 7 4 4
    4 5 2 6 5
\end{verbatim}
The maximum sum path is $7 \to 3 \to 8 \to 7 \to 5 = 30$.

\medskip
\noindent\textbf{Constraints:} $1 \le n \le 100$, each element is between 0 and 99.

\section{Solution Approach}

This is a classic dynamic programming problem. We define:
\[
  \text{dp}[i][j] = \text{maximum sum achievable from element } (i,j) \text{ down to the bottom row}
\]

\subsection{Base Case}
For the bottom row $i = n-1$:
\[
  \text{dp}[n-1][j] = \text{val}[n-1][j]
\]

\subsection{Recurrence}
For rows $i = n-2$ down to $0$:
\[
  \text{dp}[i][j] = \text{val}[i][j] + \max\!\big(\text{dp}[i+1][j],\;\text{dp}[i+1][j+1]\big)
\]

\subsection{Answer}
The answer is $\text{dp}[0][0]$.

\subsection{Why Bottom-Up?}
Bottom-up DP avoids recursion overhead and is straightforward: we start from the last
row and propagate maxima upward. The triangle can be modified in place, requiring no
extra memory beyond the input array.

\section{C++ Solution}

\begin{lstlisting}
#include <cstdio>
#include <algorithm>
using namespace std;

int tri[101][101];

int main() {
    int n;
    scanf("%d", &n);

    for (int i = 0; i < n; i++)
        for (int j = 0; j <= i; j++)
            scanf("%d", &tri[i][j]);

    // Bottom-up DP: modify triangle in place
    for (int i = n - 2; i >= 0; i--)
        for (int j = 0; j <= i; j++)
            tri[i][j] += max(tri[i + 1][j], tri[i + 1][j + 1]);

    printf("%d\n", tri[0][0]);
    return 0;
}
\end{lstlisting}

\section{Complexity Analysis}

\begin{itemize}
  \item \textbf{Time complexity:} $O(n^2)$. We visit each element of the triangle exactly once.
        The total number of elements is $\frac{n(n+1)}{2}$, so the work is $\Theta(n^2)$.
  \item \textbf{Space complexity:} $O(n^2)$ for storing the triangle. Since we modify the
        triangle in place, no additional space is needed beyond the input.
        Alternatively, one could use a 1D array of size $n$ for a rolling approach,
        yielding $O(n)$ extra space.
\end{itemize}

\end{document}
